\section{Discussion}\label{sec:discussion}

\subsection{The answer}

On holiday, do Spaniards go back to the villages their families left? Yes, in a precise sense. Cities send more of their residents to the provinces where more of them were born on every night of the year, holiday or not. What the summer adds is not a stronger pull overall, which the data cannot confirm, but a different one: the pull of residents who arrived before 1981 and of those born in villages grows, and all of the growth goes to inland destinations. At the same time, the most emptied villages, which send people away on an ordinary night, receive more than they send on almost every holiday night. The beaches take a large part of the summer, and they shape the network as a whole; the inland part of the summer follows the map of the exodus.

\subsection{What was expected and what was not}

I expected the summer to be about the coast, and in one sense it is: the change in the shape of the network, cities that send and coastal towns that receive, is the work of the beaches, and most of it follows from the new degrees. I did not expect the roots to be so visible on ordinary nights, nor the average holiday increase to be so small. The result that surprised me most is that the average hides two opposite movements. The exodus generation pulls harder in August and recent migrants pull less, as if they spent the summer somewhere other than the provinces they came from. The beach does not erase the old roots.

The two holidays also bring people back in different ways. In August the emptied villages become receivers mainly because their own residents stay at home while people from the cities arrive; the stock of second homes accounts for about half of the summer link between emptying and visitors. At Christmas the number of visitors per resident rises as well. A single holiday average would have mixed the two.

\subsection{What is new}

As far as I know, this is the first comparison of the network of where a country sleeps on holiday with the network of where its people were born, at a scale where the structure of the network can be measured. Earlier work compared holiday travel with migration between 31 provinces without controls for distance \citep{li2016spring}, or measured return tourism in one Spanish region \citep{garciagonzalez2024retorno}. The result also has a practical side: the villages that lost most people are, on holiday nights, the ones that receive more people than they send, and those visitors appear neither in the population register nor in hotel statistics.

\subsection{Limitations}

\begin{itemize}
\item \textbf{The data are estimates.} The overnight counts come from one mobile-phone operator, scaled up to the population, and each phone's home zone is inferred from where it usually sleeps. People who spend long summers in the village may be counted as residents there.
\item \textbf{The test is about cities, not people.} The roots network gives the birthplaces of a city's residents, not of the people who travel. A city with more residents born in a province sends more people there, but I cannot tell whether the travellers are those residents, their children or their neighbours. The census records the year of arrival in the province, not in the city, and the birthplace by province, not by village.
\item \textbf{Zones group small municipalities.} More than a third of the 2,435 zones (907) group several municipalities, so a ``village'' is often a group of villages. The village results hold in zones made of a single municipality, but the $k$-core result showed how much the grouping can matter.
\item \textbf{Six months of one year.} The data cover February, March, August, October, November and December 2024, with checks on 2022 and 2023. Single nights are compared, not whole seasons, and causes are not identified.
\item \textbf{Some sizes depend on choices.} The sign of the generation effects is stable, but their size depends on how city--province cells with no residents are treated (Table~\ref{tab:zero}).
\end{itemize}

\subsection{Lessons for future work}

Three choices made the difference here, and I would make them again. First, single nights: an average over a holiday period mixes the people who leave with those who come back. Second, degree-preserving null models before any interpretation: the denser holiday networks changed almost every property, and most of those changes turned out to be forced by the degrees. The rich club needs particular care: its raw value is the result, and a ratio to rewired networks close to one does not mean there is no club \citep{zhou2007structural}. Third, splitting the roots by group: the average told almost nothing, and the groups told the story. The natural next step is to replace provinces by municipalities of birth, which the census and the 2021 survey of households allow \citep{ine2023ecepov}, and to follow the same nights over several years to see whether the pull of the exodus fades as its generation ages.
